\chapter{Aceleración SIMD}
\label{ch:simd}

El trabajo inner del matcher DFS está dominado por el filtrado
de aristas: chequear \code{rel\_type\_tag},
\code{entity\_type\_tags[dst]}, inclusión de timestamp y
evaluación de predicados para millones de aristas candidatas por
hunt. Tres de esos cuatro chequeos son comparaciones a nivel byte
que vectorizan limpiamente. GraphHunter tiene dos caminos SIMD:

\begin{enumerate}
    \item \textbf{libgraphmatch} --- una biblioteca C++ enlazada
        por FFI cuando se compila el feature \code{simd} y el CPU
        host reporta soporte AVX2. El throughput peak es
        aproximadamente 3$\times$ el del matcher Rust escalar.
    \item \textbf{\code{simd\_rust} puro-Rust} (P2-D) --- una
        implementación \code{std::arch::x86\_64} que usa
        intrínsecos AVX2 sin salir del toolchain Rust. Hoy
        implementa una sola primitiva,
        \textsc{Intersect-Sorted-U32}, con más en camino.
\end{enumerate}

\section{Puente FFI de libgraphmatch}

\begin{lstlisting}[style=rust]
// graph_hunter_core/src/simd_matcher.rs (abreviado)
unsafe extern "C" {
    fn gm_graph_new() -> *mut c_void;
    fn gm_graph_add_node(g: *mut c_void, id: u32, et: u8);
    fn gm_graph_add_edge(g: *mut c_void, src: u32, dst: u32, rt: u8, ts: i64);
    fn gm_graph_finalize(g: *mut c_void);
    fn gm_matcher_run(m: *mut c_void, max_results: u32, min_score: f64)
         -> *mut c_void;
    fn gm_results_free(r: *mut c_void);
    fn gm_graph_free(g: *mut c_void);
}
\end{lstlisting}

El lado Rust envuelve los handles \code{*mut c\_void} en
\code{GmGraph} (RAII, drop llama a \code{gm\_graph\_free}) y
cachea el grafo resultante junto con un
\code{HashMap<String, u32>} que mapea string IDs a los índices
de nodo C++ (ver \code{CachedSimdGraph} en
\filepath{simd\_matcher.rs}). El caché sobrevive a los hunts en
la misma sesión; se invalida al mutar el grafo.

\textsc{Search-Temporal-Pattern-Simd} es entonces una función
corta:

\begin{algorithm}
\caption{\textsc{Search-Temporal-Pattern-Simd}}\label{alg:search-simd}
\KwIn{grafo $G$, hipótesis $H$, cap $\mathrm{max}$}
\If{AVX2 no detectado en runtime or $H$ tiene predicados}{
    \Return \textsc{Search-Temporal-Pattern-Smart}($G, H, \ldots$)
    \Comment*[r]{fallback}
}
$C \gets $ \code{G.cached\_simd\_graph()}
\Comment*[r]{construir una vez, reusar}
$r \gets $ \code{gm\_matcher\_run($C$, max, 0.0)}\;
\Return \code{translate\_node\_ids($r$, $C$.\text{string\_map})}\;
\end{algorithm}

El fallback al matcher smart en Rust es inevitable porque
libgraphmatch no evalúa predicados de metadata
(\code{\{key="val"\}}) --- éstos requieren acceso al
\code{MetadataStore}, que el lado C++ no puede hacer sin enviar
los bytes crudos por FFI.

\section{Intersección AVX2 puro-Rust}

El refactor P2-D empezó a reemplazar libgraphmatch una primitiva
a la vez. La primera que se entregó es
\code{simd\_rust::intersect\_sorted\_u32\_count}: dados dos
slices \code{\&[u32]} ascendentes deduplicados, retorna la
cantidad de elementos en común. Esta es la primitiva que
libgraphmatch usa para intersectar vecindades durante la
expansión SIMD del patrón, y es la primitiva más barata de probar
segura en aislamiento.

\begin{algorithm}
\caption{\textsc{Intersect-Sorted-U32-Count}}\label{alg:intersect}
\KwIn{slices ordenados $A$, $B$ de \code{u32}, ambos ascendentes
      y deduplicados}
\KwOut{cantidad de elementos comunes}
\If{$|A| < 32$ or $|B| < 32$}{\Return \textsc{Scalar-Merge}($A, B$)\;}
\If{AVX2 no detectado en runtime}{\Return \textsc{Scalar-Merge}($A, B$)\;}
\tcp{lado corto es $S$, lado largo es $L$}
\lIf{$|A| \leq |B|$}{$(S, L) \gets (A, B)$}\lElse{$(S, L) \gets (B, A)$}
$j \gets 0$; $\mathrm{count} \gets 0$\;
\ForEach{$s \in S$}{
    \While{$j + 8 \leq |L|$}{
        $\mathrm{chunk} \gets $ AVX2 load $L[j..j+8]$\;
        $\mathrm{needle} \gets $ AVX2 broadcast $s$ a 8 lanes\;
        $\mathrm{cmp} \gets $ \code{\_mm256\_cmpeq\_epi32}(chunk, needle)\;
        $\mathrm{mask} \gets $ \code{\_mm256\_movemask\_epi8}(cmp)\;
        \If{$\mathrm{mask} \neq 0$}{
            $\mathrm{count} \gets \mathrm{count} + 1$;
            \Break loop interno\;
        }
        \tcp{avanzar $j$ si todos los lanes del chunk son $< s$}
        \If{$L[j+7] < s$}{$j \gets j + 8$; \Continue\;}
        \Break loop interno\;
    }
    \textsc{Scan-Tail}($L, j, s, \mathrm{count}$)
    \Comment*[r]{scan escalar para $< 8$ restantes}
}
\Return $\mathrm{count}$\;
\end{algorithm}

El algoritmo es el esquema clásico ``broadcast-and-compare''
usado en la literatura de intersección SIMD. Su corrección se
desprende del invariante de que $j$ sólo incrementa
monotónicamente después de que la ventana completa $L[j..j+8]$
fue descartada (todos los elementos estrictamente menores que el
$s$ actual); como $S$ y $L$ están ordenados, ningún $s$ posterior
puede matchear en ningún índice $< j$.

\begin{theorem}[Complejidad de la intersección]
Sobre inputs de tamaños $|A|$ y $|B|$ con $|S| = \min(|A|, |B|)$,
$|L| = \max(|A|, |B|)$:
\begin{itemize}
    \item Tiempo: $\Ord{|S| \cdot (|L| / 8 + 1)}$ en el peor
        caso. El factor $/8$ es el ancho de lane AVX2. En la
        práctica la iteración por el lado corto amortiza cerca
        de $\Ord{|S| + |L| / 8}$ porque $j$ avanza
        monotónicamente.
    \item Espacio: $\Ord{1}$ auxiliar (el estado de registros).
\end{itemize}
\end{theorem}

\begin{theorem}[Seguridad en runtime]
\textsc{Intersect-Sorted-U32-Count} no invoca comportamiento
indefinido aun cuando se compile sin AVX2, porque el cuerpo
unsafe del intrínseco sólo se llama a través de
\code{is\_x86\_feature\_detected!("avx2")} en runtime. El gate
\code{\#[target\_feature(enable = "avx2")]} en compile time sólo
habilita el cuerpo del intrínseco dentro de una función que el
path escalar no llama.
\end{theorem}

\begin{proof}[Demostración (esbozo)]
El wrapper público es safe. Chequea
\code{is\_x86\_feature\_detected!("avx2")} primero; sólo en hit
invoca al cuerpo unsafe. Cuando AVX2 no se reporta, se llama al
merge escalar y no se ejecutan instrucciones AVX2. El gate
\code{target\_feature} significa que el cuerpo unsafe mismo se
compila con codegen AVX2 de todos modos, pero eso está bien
porque sólo se llama cuando el CPU lo soporta. \qed
\end{proof}

\section{Roadmap}

Primitivas restantes por portar de libgraphmatch a
\code{simd\_rust}:
\begin{itemize}
    \item \code{count\_edges\_by\_rel\_type}: barrer un
        \code{\&[CompactRelation]} buscando entradas que
        matcheen un \code{rel\_type\_tag} dado. Trivial
        ($\texttt{\_mm256\_cmpeq\_epi8}$ sobre el array de tags
        empaquetado).
    \item \code{timestamp\_range\_filter}: marcar aristas cuyo
        \code{timestamp} cae en $[t_0, t_1]$. Requiere
        comparaciones de 64 bits --- una instrucción extra por
        lane.
    \item \code{intersect\_sorted\_u32\_out}: la variante que
        materializa output de la primitiva actual sólo-count.
\end{itemize}

Cuando aterricen las tres, el path libgraphmatch de
\textsc{Search-Temporal-Pattern-Simd} puede reemplazarse por un
despacho puro-Rust, retirando la dependencia C++.

\begin{inthecode}
\filepath{graph\_hunter\_core/src/simd\_matcher.rs} --- puente
FFI.\\
\filepath{graph\_hunter\_core/src/simd\_rust.rs} ---
\textsc{Intersect-Sorted-U32-Count}.\\
\filepath{libgraphmatch/src/simd\_intersect.cpp} --- implementación
C++ de referencia.\\
\filepath{graph\_hunter\_core/src/graph.rs:689} ---
\textsc{Search-Temporal-Pattern-Simd}.
\end{inthecode}
